\begin{lemma}
  The sequence $(Z_n)$ defined by Algorithm~\ref{alg:alpha} is tight.
\end{lemma}

\begin{theorem}[Prokhorov]
  A sequence of random variables is tight if and only if every of its subsequence contains a further subsequence that converges in distribution.
\end{theorem}

\begin{theorem}[{\cite[Theorem 2.6]{billingsley1999convergence}}]
  A sequence of random variables converges in distribution if and only if every of its subsequence has a further subsequence that converges in distribution to the same limit.
\end{theorem}

With these two theorems, we have to show that if $Z_1 \sim \mu$ and $Z_2 \sim \nu$ are solutions to the fixed-point equation \eqref{eq:fixed-point}, then $\mu = \nu$.

Hence, for every $Z_1\sim \mu$ and $Z_2\sim \nu$, we have
\begin{align*}
  d_2^2(\mu, \nu)
   & \le \EE\left[(Z_1-Z_2)^2\right]                                                                                                                                                                  \\
   & = u\EE\left[\left(U^{\alpha} \left(Z_1^{(1)}-Z_2^{(1)}\right) + \left(1-U\right)^{\alpha} \left(Z_1^{(2)}-Z_2^{(2)}\right)\right)^2\right]                                                       \\
   & \,\,\,\, + v\EE\left[\left(\max\left(U^{\alpha} Z_1^{(1)}, \left(1-U\right)^{\alpha} Z_1^{(2)}\right)-\max\left(U^{\alpha} Z_2^{(1)}, \left(1-U\right)^{\alpha} Z_2^{(2)}\right)\right)^2\right] \\
\end{align*}
Using $(\max(a,b)-\max(c,d))^2 \le (a-c)^2 + (b-d)^2$ for every $a,b,c,d\in\RR$, we get
\begin{align*}
  d_2^2(\mu, \nu)
   & \le (u+v)\EE\left[\left(U^{\alpha} \left(Z_1^{(1)}-Z_2^{(1)}\right) + \left(1-U\right)^{\alpha} \left(Z_1^{(2)}-Z_2^{(2)}\right)\right)^2\right] \\
\end{align*}
Using $\EE[Z_1]=\EE[Z_2]=1$, we get
\begin{align*}
  d_2^2(\mu, \nu)
   & \le (u+v)\EE\left[U^{2\alpha} \left(Z_1^{(1)}-Z_2^{(1)}\right)^2 + \left(1-U\right)^{2\alpha} \left(Z_1^{(2)}-Z_2^{(2)}\right)^2\right] \\
   & = \dfrac{2(u+v)}{2\alpha+1}\EE\left[(Z_1-Z_2)^2\right]
\end{align*}
Taking the infimum over $Z_1\sim \mu$ and $Z_2\sim \nu$, we get
$$d_2^2(\mu, \nu) \le \dfrac{2(u+v)}{2\alpha+1}d_2^2(\mu, \nu).$$

Since $\dfrac{2(u+v)}{2\alpha+1} < 1$, we have $d_2^2(\mu, \nu) = 0$ and $\mu = \nu$. Therefore, $Z_n\xrightarrow{d} Z$, where $Z \in D_2(1)$ is the solution to the fixed-point equation \eqref{eq:fixed-point} with $\alpha = \alpha$.

Finally, we prove that $X_n\xrightarrow{d} Z$. We follow \cite{neininger2004general} to introduce an accompanying sequence $(Q_n)$ defined by
\begin{equation}
  \label{eq:accompanying-sequence}
  Q_{n+2}
  \stackrel{d}{=} I\left(A_n Z^{(1)} + B_n Z^{(2)}\right) +  J\max\left(A_n Z^{(1)}, B_n Z^{(2)}\right), n\ge 1, \quad Z_1=1, Z_2=0,
\end{equation}
where $A_n$ and $B_n$ are as in \eqref{eq:normalized-recurrence}. Let $X_n\sim \mu_n$ and $Z\sim \mu$. Let $\delta>0$. We have
$$d_2^2(\mu_{n+2}, \mu) \le \EE\left[(X_{n+2}-Z)^2\right] \le (1+\delta)\EE\left[(X_{n+2}-Q_{n+2})^2\right] + \left(1+\dfrac{1}{\delta}\right)\EE\left[(Q_{n+2}-Z)^2\right]$$

First, we show that $\EE\left[(Q_{n+2}-Z)^2\right] \to 0$. We have
\begin{align*}
  \EE\left[(Q_{n+2}-Z)^2\right]
   & \le (u+v)\EE\left[\left(A_n^{\alpha} - U^{\alpha}\right)^2 + \left(B_n^{\alpha} - \left(1-U\right)^{\alpha}\right)^2\right]\EE[Z^2] \\
   & \,\,\,\,\, + 2u\EE\left[\left(A_n^{\alpha} - U^{\alpha}\right)\left(B_n^{\alpha} - \left(1-U\right)^{\alpha}\right)\right]\EE[Z]^2  \\
   & \to 0.
\end{align*}

Next, we have
\begin{align*}
   & \EE\left[(X_{n+2}-Q_{n+2})^2\right]                                                                                                  \\
   & \le  (u+v)\EE\left[A_n^{2\alpha}\left(X_{U_n}^{(1)}-Z^{(1)}\right)^2 + B_n^{2\alpha}\left(X_{n+1-U_n}^{(2)}-Z^{(2)}\right)^2 \right] \\
   & + 2u\EE\left[A_n^{\alpha}B_n^{\alpha}\left(X_{U_n}^{(1)}-Z^{(1)}\right)\left(X_{n+1-U_n}^{(2)}-Z^{(2)}\right)\right]                 \\
\end{align*}


\section{Discussion}

\section{Appendix}
\begin{definition}[{locally convex space, \cite[Appendix 40]{zeidler1986nonlinear}}]
  \label{def:locally-convex-space}
  Let $X$ be a vector space over $\KK$ and $\{p_j\}_{j\in J}$ be a family of seminorms on $X$. Then $\left(X, \{p_j\}_{j\in J}\right)$ is called a locally convex space if for $x\in X$,
  $$x = 0 \iff p_j(x)=0 \text{ for all } j\in J.$$
\end{definition}

Let $X=C_b(S)$ where $C_b(S)$ is the space of bounded continuous functions on a metric space $S=\RR$. Equip $X$ with the supremum norm
$$\|f\|_\infty = \sup_{x\in S}|f(x)|.$$
Then $X$ is Banach. As a consequence of Riesz-Markov-Kakutani representation theorem, every $f\in X$ corresponds to a unique signed Borel measure $\mu$ on $S$ such that
$$\mu(f) = \int_S f d\mu$$
is linear and continuous. The signed Borel measures form the dual space $X^*$ of $X$.

We actually focus on $X^*$. The family of seminorms as in Definition~\ref{def:locally-convex-space} is $\{p_f\}_{f\in X}$ defined by
$$p_f(\mu) = \left|\int_S f(x) d\mu\right|.$$


\begin{definition}[Tightness]
  \label{def:tight}
  A sequence of probability measures $(\mu_n)$ on $\RR$ is tight if for every $\varepsilon > 0$, there exists $M > 0$ such that $\sup_n \mu_n(\{x: |x| > M\}) < \varepsilon$.

  A sequence of random variables $(Z_n)$ is tight if for every $\varepsilon > 0$, there exists $M > 0$ such that $\sup_n\PP(|Z_n| > M) < \varepsilon$.
\end{definition}

% Our first goal is to escape from the max operator in \eqref{eq:lower-bound-tighter}. A challenge is that $(z_n)$ fluctuates at the first iterations, and the time where it begins to increase depends on the parameters $u$ and $v$. We will show that $(z_n)$ is eventually increasing, that is,
% $$\exists N' \in \mathbb{N}, \forall n \ge N', \quad z_{n+1} \ge z_n.$$
% Together with unboundedness of $(z_n)$, if this statement is true, we have the following implications
% $$\exists N' \in \mathbb{N}, \forall n \ge N', \quad z_{n+1} \ge \max_{1 \le i \le n} z_i.$$
% Taking $N = 2N'-1$, we derive that
% $$\exists N \in \mathbb{N}, \forall n \ge N, \forall (N+1)/2\le i \le N, \quad z_{i} \ge z_{n+1-i}.$$

The next open question is then if there exists a measure of chaos that can capture whether $T_\alpha$ increases or decreases the mean, for a given $\alpha$. In particular, we consider the class of $\Phi$-entropies \cite{chafai2004entropies} that generalized both the variance and the Shannon entropy. The $\Phi$-entropy of a random variable $X$ is defined as
\begin{equation}
  \textbf{Ent}_{\Phi}(X) = \mathbb{E}[\Phi(X)] - \Phi(\mathbb{E}[X]),
\end{equation}
where $\Phi$ is a convex function. The variance is a special case of $\Phi$-entropy with $\Phi(x) = x^2$. The Shannon entropy is another special case with $\Phi(x) = x\ln x$. We leave the following question open.

\begin{openquestion}
  \label{openq:entropy}
  Is there a convex function $\Phi$ such that there exists a monotone function $\alpha: \RR_+ \to [0,1]$ such that for every $\mu$ in some subset of $\D$ and $X\sim \mu$, we have $\alpha(\textbf{Ent}_{\Phi}(X)) = \EE[X]$?
\end{openquestion}

That means we ask if it is possible to find which $\alpha$ that keeps the mean of $\mu$, just by looking at the $\Phi$-entropy of $X\sim\mu$. A more tractable question is stated and proved as follows.

statement is true and provable: for each $\mu$, there exists a unique $\alpha$ such that $\EE[T_\alpha(\mu)] = \EE[\mu]$.
